\documentclass[a4paper,12pt]{article}
\usepackage[T2A]{fontenc}
\usepackage[utf8]{inputenc}
\usepackage[russian]{babel}
\usepackage{amsmath,amssymb,mathtools}
\usepackage{helvet}
\renewcommand{\sfdefault}{cmss}
\renewcommand{\familydefault}{\sfdefault}
\usepackage{icomma}
\usepackage{geometry}
\usepackage{microtype}
\usepackage{tikz}
\usetikzlibrary{calc,arrows.meta,positioning,shapes.geometric}
\usepackage{tabularx,booktabs,longtable}
\usepackage{enumitem}
\usepackage{parskip}
\usepackage{fancyhdr}
\usepackage{setspace}
\usepackage{xcolor}
\usepackage[hidelinks]{hyperref}

\geometry{left=19mm,right=19mm,top=21mm,bottom=20mm}
\setlength{\headheight}{25pt}
\onehalfspacing
\setlength{\parindent}{0pt}
\setlist[itemize]{leftmargin=6mm,itemsep=2pt,topsep=3pt}
\setlist[enumerate]{leftmargin=8mm,itemsep=5pt,topsep=4pt}
\definecolor{accent}{HTML}{2A6F6B}
\definecolor{darkaccent}{HTML}{174B48}
\definecolor{textgray}{HTML}{2F3437}
\definecolor{softgray}{HTML}{DCE8E6}
\color{textgray}

\pagestyle{fancy}
\fancyhf{}
\fancyhead[L]{\small Кредитные задачи: три схемы}
\fancyhead[R]{\small 09.10.26}
\fancyfoot[C]{\small \thepage}
\renewcommand{\headrulewidth}{0.4pt}
\renewcommand{\headrule}{\hbox to\headwidth{\color{softgray}\leaders\hrule height \headrulewidth\hfill}}

\newcommand{\key}[1]{\textcolor{darkaccent}{\textbf{#1}}}
\newcommand{\fboxeq}[1]{\[\boxed{#1}\]}
\newcommand{\ru}{\,\text{руб.}}

\begin{document}
\thispagestyle{empty}
\begin{center}
\vspace*{31mm}
{\Large\textcolor{accent}{\textbf{Чек-лист}}}\\[9mm]
{\Huge\textbf{Кредитные задачи}}\\[6mm]
{\Large\textbf{Три схемы погашения}}\\[6mm]
{\large Аннуитетная, дифференцированная, произвольная}\\[19mm]
{\large 09.10.26}\\[31mm]
{\large Лёвин Артём Александрович}\\[3mm]
{\large\textcolor{darkaccent}{эксклюзивно для Екатерины Гнедковой}}
\end{center}
\vfill
\begin{center}
\begin{tikzpicture}[x=1cm,y=1cm]
\draw[accent!60,line width=1.15pt]
(-7,0) .. controls (-4.1,0.95) and (-1.3,-0.75) .. (1.45,0.18)
.. controls (3.65,0.85) and (5.25,-0.5) .. (7,0.25);
\end{tikzpicture}
\end{center}
\clearpage

\section*{1. Общая модель: что происходит с долгом}

В начале периода у заёмщика есть \key{остаток долга} $B$. Банк начисляет на него проценты. Затем заёмщик вносит платёж $P$; сумма, оставшаяся после платежа, переходит на следующий год.

Пусть $S$ --- первоначальная сумма кредита, $p$ --- годовая ставка в процентах, $n$ --- число лет. Введём
\[
 r=\frac{p}{100},\qquad k=1+r.
\]

\begin{center}
\renewcommand{\arraystretch}{1.24}
\begin{tabularx}{\linewidth}{@{}>{\raggedright\arraybackslash}X >{\raggedright\arraybackslash}X@{}}
\toprule
Этап & Формула \\
\midrule
Долг в начале года & $B$ \\
Начисленные проценты & $Br$ \\
Общая задолженность & $B+Br=Bk$ \\
Платёж & $P$ \\
Остаток после платежа & $Bk-P$ \\
\bottomrule
\end{tabularx}
\end{center}

\key{Главное соотношение:}
\fboxeq{B_{j}=kB_{j-1}-P_j,\quad B_0=S.}
Индекс $j$ обозначает номер года. Полное погашение за $n$ лет означает $B_n=0$.

\textbf{Связь с вкладами.} Во вкладе взнос увеличивает сумму; при кредите платёж уменьшает задолженность. Сначала необходимо установить \key{порядок операций} по условию. В рассматриваемых ниже схемах проценты начисляются за год перед очередным платежом.

\section*{2. Как определить схему по условию}

\begin{center}
\renewcommand{\arraystretch}{1.35}
\begin{tabularx}{\linewidth}{@{}>{\raggedright\arraybackslash}p{48mm} >{\raggedright\arraybackslash}X@{}}
\toprule
Схема & Сигнал в условии \\
\midrule
\key{Аннуитетная} & «Каждый год вносили \textbf{одинаковую сумму}», «равные ежегодные платежи». \\
\key{Дифференцированная} & «\textbf{Тело кредита} уменьшается на одну и ту же сумму», «основной долг погашается равными долями». \\
\key{Произвольная} & Платежи задаются отдельно либо по особому правилу; равенство платежей или долей основного долга не следует из условия. \\
\bottomrule
\end{tabularx}
\end{center}

\key{Важное различие:} равные платежи и равные части основного долга --- разные условия.

\clearpage
\section*{3. Аннуитетная схема: платёж неизменен}

Платёж каждого года равен $X$. Меняется соотношение между процентами и погашением основного долга; \key{сам платёж остаётся одинаковым}.

Для трёхлетнего кредита:
\[
\begin{aligned}
 B_1&=Sk-X,\\
 B_2&=(Sk-X)k-X=Sk^2-Xk-X,\\
 B_3&=(Sk^2-Xk-X)k-X\\
 &=Sk^3-X(k^2+k+1).
\end{aligned}
\]

Кредит полностью погашен, значит $B_3=0$. Поэтому
\fboxeq{Sk^3=X(k^2+k+1),\qquad S=\frac{X(k^2+k+1)}{k^3}.}
Если неизвестен платёж, выразите $X$ из того же равенства.

\textbf{Разобранный тип задачи.} Банк выдал кредит на три года под $10\%$ годовых. В конце каждого года вносили по $2662$ руб. Найдите первоначальную сумму.

\[
 k=1{,}1,\qquad S=\frac{2662(1{,}21+1{,}1+1)}{1{,}331}
 =\frac{2662\cdot3{,}31}{1{,}331}=6620\ru.
\]

\begin{center}
\renewcommand{\arraystretch}{1.23}
\small
\begin{tabularx}{\linewidth}{@{}c *{5}{>{\centering\arraybackslash}X}@{}}
\toprule
Год & Долг в начале & Проценты & Всего должен & Платёж & Остаток \\
\midrule
1 & 6620 & 662 & 7282 & 2662 & 4620 \\
2 & 4620 & 462 & 5082 & 2662 & 2420 \\
3 & 2420 & 242 & 2662 & 2662 & 0 \\
\bottomrule
\end{tabularx}
\end{center}
\textit{Все суммы в таблице указаны в рублях.}

\key{Как проверять решение:} после последнего платежа остаток должен равняться нулю. При вынесении $X$ в выражении $Xk^2+Xk+X$ не теряйте \textbf{последнюю единицу}.

\section*{4. Краткое повторение: вклад с пополнениями}

Для вклада $100\,000$ руб. под $10\%$ на три года с взносами $D$ в начале второго и третьего годов итоговая сумма равна
$100\,000\cdot1{,}1^3+D(1{,}1^2+1{,}1)$. Если она составляет $156\,200$ руб., то $133\,100+2{,}31D=156\,200$, откуда $D=10\,000$ руб. Взнос \key{увеличивает} сумму вклада, а платёж по кредиту \key{уменьшает} долг.

\clearpage
\section*{5. Дифференцированная схема: равные части долга}

Основной долг ежегодно уменьшается на \key{одну и ту же сумму} $S/n$. Проценты начисляют на остаток в начале года. Поэтому платежи постепенно уменьшаются при положительной ставке.

\[
 \underbrace{P_j}_{\text{платёж за год }j}
 =\underbrace{\frac Sn}_{\text{основной долг}}
 +\underbrace{\frac{S(n-j+1)}{n}\,r}_{\text{проценты}},
 \qquad j=1,\ldots,n.
\]

Отсюда удобно сразу записать первый и последний платежи:
\fboxeq{P_1=\frac Sn+Sr,\qquad P_n=\frac Sn+\frac Sn\,r.}

\textbf{Разобранный тип задачи.} Кредит $120$ тыс. руб. взят на три года под $20\%$ годовых. Каждый год погашается $40$ тыс. руб. основного долга.

\begin{center}
\renewcommand{\arraystretch}{1.23}
\small
\begin{tabularx}{\linewidth}{@{}c *{5}{>{\centering\arraybackslash}X}@{}}
\toprule
Год & Долг в начале & Проценты & Всего должен & Платёж & Остаток \\
\midrule
1 & 120 & 24 & 144 & 64 & 80 \\
2 & 80 & 16 & 96 & 56 & 40 \\
3 & 40 & 8 & 48 & 48 & 0 \\
\bottomrule
\end{tabularx}
\end{center}
\textit{Все суммы в таблице указаны в тысячах рублей.}

\key{Обратите внимание:} в колонке «платёж» каждый раз стоит $S/n$ \textbf{плюс проценты от текущего остатка}. Делить текущий остаток на $n$ каждый год нельзя.

\section*{6. Ставка по последнему платежу}

\textbf{Задача из занятия.} Кредит $7$ млн руб. рассчитан на $10$ лет. Долг погашается равными долями основного долга. Последний платёж составляет \key{не менее} $0{,}819$ млн руб. Найдите наименьшую годовую ставку $p$.

Десятая часть кредита равна $0{,}7$ млн руб. В начале последнего года осталось столько же. Значит,
\[
 0{,}7+0{,}7r\ge0{,}819
 \quad\Longleftrightarrow\quad
 0{,}7r\ge0{,}119
 \quad\Longleftrightarrow\quad r\ge0{,}17.
\]
Следовательно, \key{наименьшая ставка --- $17\%$ годовых}.

\textit{Языковая подсказка:} «не менее» $\Rightarrow\ge$, «не более» $\Rightarrow\le$.

\clearpage
\section*{7. Произвольная схема платежей}

Если размеры платежей различны и про равномерное погашение основного долга ничего не сказано, используйте общую модель \key{год за годом}.

Для двух лет с платежами $P_1$ и $P_2$:
\[
 B_1=Sk-P_1,\qquad B_2=(Sk-P_1)k-P_2.
\]
При полном погашении:
\fboxeq{Sk^2=P_1k+P_2.}

\key{Метод:} составьте таблицу с колонками «долг в начале», «проценты», «общая задолженность», «платёж», «остаток». Не приписывайте задаче аннуитетную или дифференцированную схему, если в условии нет нужных признаков.

\section*{8. На что особенно смотреть в решении}

\begin{itemize}
\item \key{Обозначения.} Назовите все величины: $S$, $p$, $r$, $k$, $P_j$ или $X$. Запись «процент $=10\%$» без пояснения величины неудачна.
\item \key{Суммы и проценты.} Если остаток $B$, начисленные проценты равны $Br$, а задолженность перед платежом --- $Bk$.
\item \key{Признаки схем.} «Одинаковый платёж» относится к аннуитету; «одинаковая часть тела кредита» --- к дифференцированному погашению.
\item \key{Скобки.} В аннуитетной формуле $X(k^2+k+1)$ последнее слагаемое $1$ обязательно.
\item \key{Неравенства.} «Не менее» и «не более» сохраняйте в математической записи с правильным знаком.
\item \key{Числовые вычисления.} Записывайте промежуточные подстановки и действия. Контролируйте единицы: рубли, тысячи, миллионы.
\item \key{Проверка.} Конечный остаток при полном погашении равен $0$; платежи не должны делать долг отрицательным раньше срока.
\end{itemize}

\clearpage
\thispagestyle{plain}
\begin{center}
{\Large\textbf{Алгоритм: как выбрать схему кредита}}
\end{center}
\vspace{8mm}
\begin{center}
\begin{tikzpicture}[
  startstop/.style={ellipse,draw=accent,line width=0.95pt,align=center,minimum width=44mm,minimum height=11mm,fill=white},
  process/.style={rectangle,rounded corners=2mm,draw=accent,line width=0.95pt,align=center,text width=64mm,minimum height=13mm,inner sep=2mm,fill=white},
  choice/.style={rectangle,rounded corners=2mm,draw=accent,line width=0.95pt,align=center,text width=39mm,minimum height=15mm,inner sep=2mm,fill=white},
  decision/.style={diamond,aspect=2.7,draw=accent,line width=0.95pt,align=center,text width=33mm,inner sep=1.5mm,fill=white},
  arr/.style={-{Stealth[length=2.8mm]},line width=1pt,draw=darkaccent},
  every node/.style={font=\small}
]
\node[startstop] (begin) at (0,0) {Начало};
\node[process] (read) at (0,-1.7) {Прочитайте условие и выделите правило погашения};
\node[decision] (rule) at (0,-4.1) {Как задано правило платежей?};
\node[choice] (annu) at (-5.25,-7.05) {Равные платежи\\аннуитетная схема};
\node[choice] (diff) at (0,-7.05) {Равные доли тела\\дифференцированная схема};
\node[choice] (other) at (5.25,-7.05) {Иное правило\\произвольная схема};
\node[process] (balance) at (0,-10.0) {Для каждого года:\\$B_j=(1+r)B_{j-1}-P_j$};
\node[startstop] (finish) at (0,-12.1) {Проверьте $B_n=0$};
\draw[arr] (begin)--(read);
\draw[arr] (read)--(rule);
\draw[arr] (rule.west) -| node[pos=0.22,above]{равные суммы} (annu.north);
\draw[arr] (rule.south) -- node[right,pos=0.60]{равные доли} (diff.north);
\draw[arr] (rule.east) -| node[pos=0.22,above]{иное правило} (other.north);
\draw[arr] (annu.south) |- (balance.west);
\draw[arr] (diff.south)--(balance.north);
\draw[arr] (other.south) |- (balance.east);
\draw[arr] (balance)--(finish);
\end{tikzpicture}
\end{center}

\clearpage
\section*{9. Тренировка: Путь А}
\textit{10 заданий. Выполните вычисления самостоятельно, записывая ключевые промежуточные шаги. В заданиях о кредитах платежи производятся в конце каждого года после начисления процентов.}

\begin{enumerate}[label=\textbf{\arabic*.},leftmargin=8mm,itemsep=8pt]
\item Укажите схему для каждого условия: (а) ежегодно вносят одну и ту же сумму; (б) каждый год погашают $1/5$ первоначального основного долга плюс проценты; (в) в первый год платят $500$ руб., во второй --- $700$ руб., иных правил нет.

\item Кредит $S=100$ тыс. руб. выдан на три года под $10\%$ годовых, платёж каждого года равен $X$. Выпишите остатки долга после первого, второго и третьего платежей. Составьте уравнение полного погашения. \textit{Решать его численно не требуется.}

\item Трёхлетний кредит под $10\%$ годовых погасили тремя равными платежами по $1331$ руб. Найдите первоначальную сумму кредита.

\item Кредит $4200$ руб. взят на два года под $10\%$ годовых. Найдите размер одинакового ежегодного платежа, если кредит полностью погашен.

\item Кредит $150$ тыс. руб. выдан на три года под $12\%$ годовых. Основной долг погашают равными долями. Составьте таблицу по годам: проценты, платёж, остаток.

\item Кредит $600$ тыс. руб. выдан на четыре года под $15\%$ годовых с погашением основного долга равными долями. Найдите первый, последний и суммарный платежи.

\item Кредит $9$ млн руб. выдан на $15$ лет с равномерным погашением основного долга. Последний платёж --- не менее $0{,}66$ млн руб. Найдите наименьшую возможную ставку в процентах годовых.

\item Кредит $800$ тыс. руб. выдан на четыре года с погашением основного долга равными долями. Первый платёж равен $280$ тыс. руб. Найдите годовую ставку и последний платёж.

\item Кредит $1000$ руб. выдан на два года под $10\%$ годовых. После первого года внесли $600$ руб. Сколько потребуется внести в конце второго года, чтобы закрыть долг? Схема платежей заранее не задана.

\item Кредит $2100$ руб. взят на два года под $10\%$ годовых. Сравните аннуитетную и дифференцированную схемы: найдите платежи каждого года по обеим схемам и укажите, при какой из них суммарная выплата меньше и на сколько.
\end{enumerate}

\clearpage
\section*{10. Ответы}
\textit{Проверьте не только итоговые числа, но и выбранную схему, начисление процентов и нулевой остаток после полного погашения.}

\renewcommand{\arraystretch}{1.38}
\begin{longtable}{@{}>{\raggedright\arraybackslash}p{15mm} >{\raggedright\arraybackslash}p{144mm}@{}}
\toprule
№ & Ответ / контрольный результат \\
\midrule
\endfirsthead
\toprule
№ & Ответ / контрольный результат \\
\midrule
\endhead
1 & (а) аннуитетная; (б) дифференцированная; (в) произвольная. \\
2 & $B_1=110-X$, $B_2=121-2{,}1X$, $B_3=133{,}1-3{,}31X$ (тыс. руб.); $133{,}1-3{,}31X=0$. Здесь $X$ выражен в тыс. руб. \\
3 & $3310$ руб. \\
4 & $2420$ руб. \\
5 & Проценты: $18$, $12$, $6$ тыс. руб.; платежи: $68$, $62$, $56$ тыс. руб.; остатки: $100$, $50$, $0$ тыс. руб. \\
6 & Первый --- $240$ тыс. руб.; последний --- $172{,}5$ тыс. руб.; всего --- $825$ тыс. руб. \\
7 & $10\%$ годовых. \\
8 & $10\%$ годовых; последний платёж $220$ тыс. руб. \\
9 & $550$ руб. \\
10 & Аннуитетная: $1210$ и $1210$ руб. (всего $2420$ руб.). Дифференцированная: $1260$ и $1155$ руб. (всего $2415$ руб.). Дифференцированная дешевле на $5$ руб. \\
\bottomrule
\end{longtable}

\vfill
{\small\textcolor{darkaccent}{\textbf{Самопроверка.}} Понимаю, откуда берётся каждое слагаемое платежа; отличаю равный платёж от равной части долга; сохраняю знак неравенства из условия; умею проверить остаток после последнего года.}
\end{document}
